\chapter{반(수강반) 관리}

\begin{summarybox}{한눈에}
학원 안에서 학생들을 \textbf{같은 시간·같은 학년·같은 강의} 단위로 묶는 ``수강반'' 을 만들고 관리합니다. 학습 계획표는 반 단위로도 발급할 수 있어, 반을 잘 짜놓으면 운영이 훨씬 빠릅니다.
\end{summarybox}

\kfShot{03-classes}{반 관리 — 본 기관의 수강반 목록}

\section*{화면 구조}

\begin{screenbox}{반 목록 + 신규 생성}
\noindent
\begin{tabular}{@{}p{3cm}p{2.5cm}p{2cm}p{2cm}p{2cm}@{}}
\toprule
\textbf{반 이름} & \textbf{레벨/학년} & \textbf{학생 수} & \textbf{상태} & \textbf{액션} \\
\midrule
초3 월수금 A & 소쉬르 3 / 초3 & 8 & active & \inacttab{학생 관리} \inacttab{수정} \\
중1 화목 B & 러셀 1 / 중1 & 12 & active & \inacttab{학생 관리} \inacttab{수정} \\
\bottomrule
\end{tabular}

\medskip
\activetab{+ 새 반 만들기}
\end{screenbox}

\section*{여기서 할 수 있는 일}
\begin{itemize}[leftmargin=1.2em]
  \item \textbf{새 반 만들기} — 이름 / 레벨(소쉬르~비트겐슈타인) / 학년 / 시작일 입력
  \item 반에 \textbf{학생 추가/제외}
  \item 반 \textbf{비활성화} (시즌 종료 시 — 데이터는 보존)
  \item 반 \textbf{이름·설명·레벨 수정}
\end{itemize}

\section*{자주 쓰는 흐름}
\begin{enumerate}
  \item 학기 시작 → ``+ 새 반 만들기'' → 이름·레벨·학년 입력 → 생성
  \item 학생 관리에서 학생들을 한 번에 선택 → 이 반에 추가
  \item 학습 계획표 발급 시 ``대상: 이 반 학생 전원'' 으로 한 번에 발급
\end{enumerate}

\begin{tipbox}{알아두면 좋은 점}
\begin{itemize}[leftmargin=1.2em]
  \item 반은 본 기관 안에서만 보임 — 다른 기관 학생을 같은 반에 넣을 수 없습니다
  \item 학생은 여러 반에 동시 소속 가능 (예: 정규반 + 시험 대비반)
  \item 비활성화한 반은 학생 화면에서 사라지지만 데이터(통계·과제)는 보존됩니다
\end{itemize}
\end{tipbox}
